\documentclass[12pt]{article}
\usepackage{fullpage}
\usepackage[shortlabels]{enumitem}
\usepackage{amsmath}
\usepackage{amsthm}
\usepackage{graphicx}
\usepackage{amssymb}
\usepackage{enumitem}
\usepackage{hyperref}

\title{250 Homework $\#$3 (Fall 2026)}
\author{Toby Hansen \footnote{Collaborated with Nobody.}}
%\author{Solution Key \footnote{Collaborated with Nobody.}}
\begin{document}
	\date{}
	\maketitle

	This assignment is due at 11:59 p.m.~on Friday 23 October 2026, to be
	submitted on Gradescope. Late assignments (with penalty, unless excused) will
	be accepted until 11:59 p.m. on Saturday 24 October 2026, and not after. Please
	see the syllabus for more details on homework grading and late penalties.

	\medskip

	When submitting to Gradescope, please do not forget to assign pages to problems.

	\medskip

	All problems, with the exception of problem 2, originally come from the
	textbook. Some of them have been modified. In case of doubt, the version printed
	here is the one you should follow.

	\bigskip

	The instructions for this homework are the same as for HW1 and HW2. Please read
	them there. Also, do not forget to fill out the questionnaire describing collaboration
	and use of LLMs that appears as a separate submission on Gradescope.

	% The main Canvas Homework page gives detailed submission guidelines, split into 6 major areas, listed below.  Please see the Canvas Homework page for the complete  guideline details.

	% \begin{enumerate}
	% \item {\bf Proper page formatting.}
	% (Use of Latex is recommended).
	% \item {\bf Proper proof writing and formatting.}\\
	% Proofs are marked on understandability, as well as correctness.
	% \item {\bf Only use notation taught in class.}
	% \item {\bf Use specified  proof structures. }
	% \item {\bf Properly declare collaborations and resources used.}\\
	% Fill out the brief individual form for each question.\\
	% This will be available as a survey on Canvas.
	% \item {\bf Proper Use of LLMs.}
	% \end{enumerate}

	% Solutions that do not follow guidelines  1-4 will not be marked.
	% Submissions  not satisfying  1 and 2 will will be returned for rewriting and resubmission. and  points will be deducted for the original solutions not following guidelines 1 and  2;   solutions not following 3 and 4 will not be marked.
	% Questions that do not have their acknowledgment form (guidelines 5 and 6) completed, will not be marked.

	\bigskip

	See the Canvas Homework web page for detailed instructions on how to format solutions.
	Solutions that do not follow those formatting instructions will be returned
	for rewriting and resubmission. In later assignments there will be points deducted
	for such formatting errors.

	\newpage
	%%%%%%%%%%%%%%%%%%%%%%%

	\section*{\textbf{Q1. Modified from P3.1.1} [5 pts]}

	The primes have a strange distribution. See, e.g., \url{https://en.wikipedia.org/wiki/List_of_prime_numbers}
	for a list of the first 1000 primes. It shows that, the primes are scattered.
	Sometimes two consecutive primes only differ by two, e.g., 7589 and 7591. Sometimes
	they differ by quite a lot and you get a long run of composite numbers. E.g.,
	7759 is prime and the next prime after it is 7789, so the run of 29 numbers starting
	with 7760 are all composite.

	For this problem find 75 consecutive naturals that are all composite. (\textbf{Hint}:
	The \textbf{factorial} of a natural $k$, written $k!$, is the product
	$1 \cdot 2 \cdot . . . \cdot k$. You can tell a lot about which numbers divide
	$k!$.) \\

	\textit{Note: To solve this problem you only need to write down the first number
	in the series and state that the next 74 numbers are all composite, explaining
	why you know this. The first item can be written using factorial notation and summation.}

	\subsection*{\textbf{Solution:}}

	Any number $k$'s factorial can be divided evenly by all positive numbers less than
	or equal to $k$, and adding any number $n$ to the factorial that is between
	$2$ and $k$ will result in a number that is still evenly divisble by $n$, which
	means it is composite. Using this logic, we can easily find $k$ consecutive composite
	naturals where $k$ is greater than $1$ using $(k+1)!+n$ where $n$ is between $2$
	and $k+1$. $k+1$ must be used instead of just $k$ because it isn't a guarantee
	that $k!+1$ or $k!+k+1$ are composite. \\

	With this in mind, to find 75 consecutive composite naturals, we just have to
	find:
	\[
		\{76!+2,76!+3,76!+4,\ldots,76!+76\}
	\]\\

	\newpage

	\section*{\textbf{Q2. New Problem from Section 3.5} [10 pts]}
	\begin{enumerate}[(a)]
		\item Use the extended gcd algorithm to find integers $a$ and $b$ satisfying
			$a \cdot 137 + b \cdot 111 = 1$. Show all steps of the algorithm. \\
			For full credit you must use the exact technique we taught in class, which
			is designed to catch errors early. Use of a different technique will lead to
			points being deducted, even for a correct final result.

			{\em \small Note: You should check the solution of (a) before proceeding to (b) and (c). Errors in (b) and (c) resulting from finding an incorrect $a,b$ pair in part (a) will not receive partial credit. }

		\item Find a multiplicative inverse of 137 modulo 111 and a multiplicative inverse
			of 111 modulo 137. The inverses you derive should be the smallest natural numbers
			that are inverses.

		\item Find the smallest natural solution to the congruences $x \equiv 2$ (mod
			137) and $x \equiv 3$ (mod 111). Show your work.
	\end{enumerate}

	\subsection*{\textbf{Solution:}}
	\begin{enumerate}[(a)]
		\item $137\mod111=26$\\
			$111\mod26=7$\\
			$26\mod7=5$\\
			$7\mod 5=2$\\
			$5\mod2=1$\\
			$2\mod1=0$\\
			\\
			$26=137(1)+111(-1)$\\
			$7=111(1)+26(-4)=111(1)+137(-4*1)+111(-4*-1)=137(-4)+111(5)$\\
			$5=26(1)+7(-3)=137(1)+111(-1)+137(-3*-4)+111(-3*5)=137(13)+111(-16)$\\
			$2=7(1)+5(-1)=137(-4)+111(5)+137(-1*13)+111(-1*-16)=137(-17)+111(21)$\\
			$1=5(1)+2(-2)=137(13)+111(-16)+137(-2*-17)+111(-2*21)=$\\
			$137(47)+111(-58)$\\
			Thus, $a=47$, and $b=-58$

		\item The inverse for $137\mod111$ is already given by $a$ ($47$). Since $b$
			($-58$) is negative, we can add $137$ to get the \textit{natural} inverse
			of $111\mod137$, which is $79$.

		\item By Chinese Remainder Theorem,
			$c=137(47)(3)+111(-58)(2)=19317-12876=6441$, removes the need to brute force
			a solution:\\
			$137(0)+2=2$ and $111(0)+3=3$\\
			$137(1)+2=139$ and $111(1)+3=114$\\
			\ldots\\
			$137(47)+2=6441$ and $111(58)+3=6441$\\
			\\
			Thus, $x\equiv6441\mod(137*111)$, so, $x\equiv6441\mod(15207)$\\
			Since $6441$ is less than $15207$, $6441$ must be the smallest natural
			solution.
	\end{enumerate}

	\newpage

	\section*{\textbf{Q3. Modified from P3.3.4} [10 pts]}
	We have defined the \textbf{factorial} $n!$ of a natural $n$ to be the product
	of all the naturals from $1$ through $n$, with $0!$ being defined as $1$. Let
	$p$ be an odd prime number.

	\begin{itemize}
		\item[(a)] Prove that the only solutions to $x^{2} \equiv 1\  (\bmod \ p)$ satisfying
			$0 < x < p$ are $x=1$ and $x=p-1.$

			{\em \small {\bf Hint:} Note that $x^{2} \equiv 1 \ (\bmod \ p)$ is equivalent to $(x^{2}-1) \equiv 0\  (\bmod \ p)$ and that $x^{2} -1 = (x-1) (x+1).$ }

		\item[(b)] Prove that $(p - 1)!$ is congruent to $-1$ modulo $p$.

			{\em \small {\bf Hint:} Use part (a) and also pair as many numbers as you can with their multiplicative inverses.

			{\bf Examples:} $4!=24$, $6!=720$ and $10!=3628800$ while $24 \equiv -1 \ (\bmod\  5),$ $720 \equiv -1 \ (\bmod \ 7),$ and $3628800 \equiv -1 \ (\bmod \ 11)$

			}
	\end{itemize}

	\subsection*{\textbf{Solution:}}

	\begin{itemize}
		\item[(a)] Since $(x-1)(x+1)\equiv0\mod p$, $x-1$ and $x+1$ must be
			multiples of $p$.\\
			First, since $0<x<p$, $x-1$ must be greater than or equal to $0$, but less
			than $p-1$; Since $p$ is greater than $2$ (on account of being an odd prime),
			the only multiple of $p$ in that range is technically $0$, $x-1=0$ so
			$x=1$.\\
			Second, since $0<x<p$, $x+1$ must be greater than $1$, but less than $p+1$;
			Since $p$ is greater than $2$ (on account of being an odd prime), the only
			multiple of $p$ in that range is $p$ itself, $x+1=p$ so $x=p-1$.

		\item[(b)] From part $(a)$, the only naturals that are their own
			multiplicative inverses $\mod p$ (between zero and p) are $1$ and $p-1$, so
			every natural from $\{2,3,\ldots,p-2\}$ can be paired with another natural
			in the set so that their product $\mod p=1$. For example, if $p=7$, then
			$2$ can be paired with $4$ since $8\mod7=1$, and $3$ can be paired with $5$
			since $15\mod7=1$.\\
			Since each pair's product $\mod p=1$, the product of all of them $\mod p$
			must also be equal to $1$; but that's just (p-2)!, so multiply it by p-1 to
			get (p-1)!, and since p-1 is its own multiplicative inverse, it can be
			paired with itself (in the earlier example $p-1=6$, $36\mod 7 = 1$).\\
			In practice, $6!=1*2*3*4*5*6$, but since $(2*4)$ and $3*5$ $\mod p=1$,
			only $1$ and $6$ remain when the modulo is taken: $6!\equiv1*1*1*6 \mod7$,
			which is the same as saying $6!\equiv6\equiv-1\mod 7$. $7$ was an arbitrary
			odd prime, but the example can generalized to all odd primes based on the
			logic of the set between $1$ and $p-1$ being uniquely pairable.
	\end{itemize}

	\newpage
	\section*{\textbf{Q4. P3.5.10 Modified} [15 pts]}
	Lecture 18 will discuss how to use {\em induction} to prove facts about a particular
	change-making problem. In this question
	%(which is very similar but  more general)
	you only need to use number-theoretic facts and not induction.

	Let $m$ and $n$ with $m < n$ be two positive naturals. We say that natural $x$
	is {\em $(m,n)$-good} if $x$ can be expressed as a linear combination $a \cdot
	m + b \cdot n$ where $a$ and $b$ are naturals, i.e., non-negative integers.
	%If $x$ is not {\em $(m,n)$-good}  it will be {\em $(m,n)$-bad.}  If $m,n$ are fixed we will just write {\em good} and {\em bad}.

	In parts (c) and (d), $[x]$ denotes the equivalence class of $x$ modulo $m.$

	Parts (a) and (b) introduce you to the setup. Parts (c)-(f) are proving the Theorem
	that ``There are an infinite number of {\em $(m,n)$-bad} naturals if and only if
	$m$ and $n$ are not relatively prime.''

	\begin{enumerate}[a)]
		\item Let $m = 2$ and $n = 3.$ Prove that all naturals except $1$ are {\em $(2,3)$-good}
			and $1$ is {\em $(2,3)$-bad}.

		\item Determine all naturals which are {\em $(3,5)$-good} and all which are {\em $(3,5)$-bad}.
			Prove your statement.
		% \item Let $x_1,x_2$ be naturals with $0 \le x_1 < x_2 < m.$ Prove that $ x_1 n \not\equiv x_2 n \, (\bmod \ m)$, i.e., that $[x_1 n] \not = [x_2 n],$ where $[i]$ is the equivalence class modulo $m.$

		\item Let $m,n$ be relatively prime and $i \in \{0,1,2,\ldots,m-1\}.$ Prove that
			there for every such $i$ there exists natural $j_{i}$ such that $j_{i} \cdot
			n \equiv i \, (\bmod \ m)$. \\
			Equivalently prove that there exists a natural $j$ satisfying
			$[j_{i} \cdot n]= [i]$.

		\item Suppose that $x$ is {\em $(m,n)$-good}. Prove that if $y \ge x$ and
			$y\in [x],$ then $y$ is {\em $(m,n)$-good}.

		\item Let $m,n$ be relatively prime. Prove that only a finite number of
			naturals are {\em $(m,n)$-bad.} Equivalently, prove that there exists natural
			$N,$ such that if $x \ge N,$ $x$ is {\em $(m,n)$-good}.

		\item Suppose $m,n$ are NOT relatively prime.\\
			Prove that the set of {\em $(m,n)$-bad} naturals is infinite.
	\end{enumerate}

	{\em \small Hints and comments: \begin{enumerate}[1.]\item $0$ is always $(m,n)$-good for all $m,n.$

	\item For (c), you may either prove from scratch or use any fact that we have proven in class.

	\item For (e), use parts (c) and (d). If you get stuck, you might want to look ahead to the example shown at the end of Lecture 18 to see how this can be helpful.

	\item For (f), construct an infinite sequence, all of whose elements are $(m,n)$-bad.\end{enumerate} }

	\subsection*{\textbf{Solution:}}
	\begin{enumerate}[(a)]
		\item Since $(m,n)$ is good only when $x$ is a combination of multiples of $m$
			and $n$, and EVERY natural except for $1$ can be represented by a
			combination of multiples of $m$ and $n$ (even zero), all naturals except $1$
			are {\em $(2,3)$-good}. To cover some base cases, $2=1(2)+0(3)$,
			$3=0(2)+1(3)$, $4=2(2)+0(3)$, $5=1(2)+1(3)$, etc; for any natural $n$ (other
			than $1$), if it's even it can be represented as a multiple of just $2$ like
			$k(2)+0(3)$ where $k=n/2$, if it's odd it can be represented using $2$ to
			get to $n-3$ then add a single $3$ like $j(2)+1(3)$ where $j=(n-3)/2$. Any
			number lower than $2$ can't be represented however, with the exception of
			zero, since $0=0(2)+0(3)$; since $1$ the only number lower than $2$ other
			than zero, it can't be represented my $m$ and $n$ and is therefore {\em $(m,n)$-bad}.

		\item To cover some base cases, $0=0(3)+0(5)$, $3=1(3)+0(5)$, $5=0(3)+1(5)$,
			$6=2(3)+0(5)$, so $\{0,3,5,6\}$ are {\em $(3,5)$-good}, and $\{1,2,4,7\}$ are
			{\em $(3,5)$-bad} since they can't be represented by $3a+5b$. Now, by adding
			$3$ or $5$ to any confirmed {\em $(3,5)$-good} natural, it can be
			represented by just increasing $a$ or $b$ by 1, and since EVERY natural greater
			than 7 can be acquired by adding multiples of 3 and five to $\{0,3,5,6\}$,
			every natural greater than 7 is {\em $(3,5)$-good}, therefore $\{1,2,4,7\}$
			are the only {\em $(3,5)$-bad} naturals.

		\item Since $m$ and $n$ are relatively prime, $\gcd(m,n)=1$, so by extended gcd,
			$a$ and $b$ must exist such that $am+bn=1$, so $bn\equiv1\mod m$, meaning
			$b$ is a multiplicative inverse of $n\mod m$. Multiplying both sides by $i$
			gives $ibn\equiv i\mod m$. We can just make $j_{i}=ib$, which proves that for
			every such $i$ there exists natural $j_{i}$ such that $j_{i} \cdot n \equiv
			i \, (\bmod \ m)$. Since $j_{i}$ is a multiple of $i$ ($ib$, to be precise),
			we can equivalently find a version with an even divisor instead, which satisfies
			$[j_{i} \cdot n]= [i]$.

		\item Since $x$ is {\em $(m,n)$-good}, there exists naturals $a$ and $b$
			such that $x=am+bn$. Since $y\geq x$ and $y\equiv x\mod m$, subtracting x
			from both sides gives $y-x\equiv 0\mod m$ so $y-x$ is evenly divisible by
			$m$. Therefore, $y-x=km$ where k is some natural. Substituting $x=am+bn$ into
			$y=x+km$ gives $y=am+bn+km=m(a+k)+bn$. Since $a+k$ and $b$ are both
			naturals, $y$ is also {\em $(m,n)$-good}.

		\item Since $m$ and $n$ are relatively prime, from part (c) we know that for
			every remainder $i$ there exists a natural $j_{i}$ such that
			$j_{i} n\equiv i\mod m$. Each $j_{i} n$ is {\em $(m,n)$-good} because $j_{i}
			n=0(m)+j_{i}(n)$. Since there are only $m$ possible remainders let $N$ be
			the largest. Now, for any natural $y\geq N$, let $i$ be the remainder of $y
			\mod m$. We know $y\equiv j_{i} n\pmod m$, and since $y\geq N\geq j_{i} n$,
			part (d) says $y$ is also {\em $(m,n)$-good}. Therefore, EVERY natural
			greater than or equal to $N$ is $(m,n)$-good, meaning there's a finite amount
			of naturals between 0 and $N$ that could be {\em $(m,n)$-bad}.

		\item Since $m$ and $n$ are not relatively prime, $d=\gcd(m,n)>1$. Since $d$
			divides both $m$ and $n$, it must also divide every combination of $am+bn$.
			Therefore, every {\em $(m,n)$-good} natural must be divisible by $d$. However,
			there are infinitely many natural numbers that are not divisible by $d$ ($f
			d+1$ where f is any natural). None of these numbers can be written as $am+b
			n$, so they are all {\em $(m,n)$-bad}. Therefore, there are infinitely
			many {\em $(m,n)$-bad} naturals when $m$ and $n$ are not relatively prime.
	\end{enumerate}

	\newpage

	\section*{\textbf{Q5. P3.6.4} [10 pts]}
	In this problem we will work with strings over the alphabet $\{a, b\}$, and
	let our “multiplication” be string concatenation.

	We say that one string $u$ “divides” another string $v$ if
	$\exists x : \exists y : xuy = v$, or equivalently, if $u$ is a substring of
	$v$. This is denoted by $D(u,v)$.

	As with the naturals, we redefine the word “prime” for strings so that $P(w)$
	means “$w \not = \lambda$ and if $D(x,w)$ is true, then either $x = \lambda$
	or $x = w$”

	\begin{enumerate}[(a)]
		\item What strings are “prime” using this definition?

		\item Explain why any string factors uniquely as a product (concatenation) of
			these “primes”.

		\item Prove a version of the Atomicity Lemma for these “primes”.

		\item Show that any nonempty string that is \textit{not} one of your
			“primes”” fails to satisfy the Atomicity Lemma, even in the simple version.
	\end{enumerate}
	{\itshape Notes: \begin{enumerate}[(i)]\item You can prove (c) by following the structure of the proof of the original atomicity lemma in the book and first proving a simple version and then a complete version.

	\item The proof of the simple version should be easier than the proof for naturals. You can write it in words. You may assume that given a character 'c' and a string 's', you are always able to tell whether 'c' is in 's' or not.\end{enumerate} }

	\subsection*{\textbf{Solution:}}
	\begin{enumerate}[(a)]
		\item All possible strings are $\{\lambda,a,b,aa,ab,ba,bb,aaa,etc...\}$ just
			like binary. The only prime strings using the given definition are just p[lain
			$a$ and $b$. Since there are an infinite number of combinations that can all
			"divide" each other, only the building blocks themselves are prime.

		\item Since $a$ and $b$ are the only prime strings, every nonempty string
			can be factored into a sequence of these primes, and the order matters.
			For example, $abba=a*b*b*a$, NOT $a*a*b*b$; even though it has the same amount
			of the prime characters, the order would yield a different unique string. Therefore,
			every nonempty string factors uniquely. Also, the nonempty string is just
			the lack of either of the primes and is still unique itself.

		\item If a prime divides a product, then it must divide at least one of the
			factors. In this case, the prime strings are only $a$ and $b$. Suppose one
			of the prime strings $p$ divides the concatenation $uv$. Since $p$
			consists of exactly one character, it cannot be a substring between $u$ and
			$v$. Therefore, that character must appear entirely independently within $u$,
			$v$, or both, meaning $D(p,u)$ or $D(p,v)$ or both.

		\item For any nonempty non-prime string $w$, it must be at least $2$
			characters long. We can split $w$ into two nonempty strings $u$ and $v$,
			where $u$ is the first character and $v$ is the remaining characters, so $w
			=uv$. $D(w,uv)$ is true because $uv=w$. However, $D(w,u)$ and $D(w,v)$ are
			both false because $w$ is longer than both, so it can't be a substring of either.
	\end{enumerate}

	\newpage
	\section*{\textbf{Q6. P4.1.6} [10 pts]}
	(uses Java) Give a recursive definition of and a recursive static method for
	the natural subtraction function, with pseudo-Java header

	\begin{center}
		\texttt{natural minus (natural $x$, natural $y$).}
	\end{center}
	On input $x$ and $y$ this function returns $x - y$ if this is a natural (i.e.,
	if $x \geq y$) and 0 otherwise.

	\subsection*{\textbf{Solution:}}

	\begin{verbatim}
// This function returns the difference between between two naturals, but only 
// when x is greater than or equal to y. Otherwise, it treats them as equal. 
// This is accomplished by repeatedly subtracting 1 from both both naturals 
// until y reaches zero, at which point x is equal to the raw difference.
natural minus(natural x, natural y) {
    if (y == 0) return x;
    if (x == 0) return 0;
    return minus(x - 1, y - 1);
}
\end{verbatim}
	If $y=0$, then the result is $x$, since nothing remains to be subtracted. If
	$x=0$ while $y>0$, then the result is $0$, since the subtraction would
	otherwise be negative.

	\newpage

	\section*{Q7 Modified 4.3.1 [10 pts]}

	Define $S_{n} = \sum_{i=0}^{n} i^{3},$ i.e., as the sum of the first $i$ cubes
	. Prove by induction that, for all naturals $n,$
	\begin{equation}
		S(n) = \frac{n^{2}(n+1)^{2}}{4}.
	\end{equation}

	Your proof should be split into three parts

	\begin{enumerate}[(i)]
		\item State your inductive hypothesis $P(n).$ \\
			(This must explicitly state what set $n$ comes from.)

		\item State and justify the correctness of the base case.

		\item Prove the correctness of the induction step. \\
			This must explicitly state exactly in which step(s) the proof is using the
			inductive hypothesis.\\
			Also, every line of your proof should follow easily from the previous line.
			It's your job to provide justification of each step from the previous one,
			not the reader's.
	\end{enumerate}

	\subsection*{\textbf{Solution:}}
	\begin{enumerate}[(i)]
		\item Assume that the statement is true for some arbitrary natural $k$,
			meaning $\sum_{i=0}^{k}i^{3}=\frac{k^{2}(k+1)^{2}}{4}$.

		\item For the base case, let $n=0$. The left side is $\sum_{i=0}^{0}i^{3}=0^{3}
			=0$, and the right side is $\frac{0^{2}(0+1)^{2}}{4}=0$. Since both sides are
			equal, the statement holds for $n=0$.

		\item For $k+1$:
			\[
				\begin{aligned}
					\sum_{i=0}^{k+1}i^{3} & =\sum_{i=0}^{k}i^{3}+(k+1)^{3}  & \text{(expand the summation)}                \\
					                      & =\frac{k^2(k+1)^2}{4}+(k+1)^{3} & \text{(inductive hypothesis)}                \\
					                      & =\frac{k^2(k+1)^2+4(k+1)^3}{4}  & \text{(multiply right term by 4/4, combine)} \\
					                      & =\frac{(k+1)^2(k^2+4(k+1))}{4}  & \text{(take out like terms)}                 \\
					                      & =\frac{(k+1)^2(k^2+4k+4)}{4}    & \text{(distribute the 4)}                    \\
					                      & =\frac{(k+1)^2(k+2)^2}{4}.      & \text{(factor)}
				\end{aligned}
			\]
	\end{enumerate}

	\newpage

	\section*{Q8 Modified 4.3.3 [10 pts]}

	Let the finite sequence $a_{0},a_{1},\ldots,a_{n}$ be defined by the rule $a_{i}
	= a_{0} \cdot r^{i}$, where $r$ is some number not equal to $1.$ Prove by
	induction on $n$ that
	\[
		\sum_{i=0}^{n} a_{i} = \frac{a_{0}- a_{n+1}}{1-r},
	\]
	where $a_{n+1}$ is defined by the same rule even though it is not part of the
	sequence.

	\medskip

	{\em \small Note: This is just the formula for the sum of a geometric sequence. But, you cannot use that fact. The problem is asking you to prove this directly from scratch using simple induction.}

	\medskip

	Your proof should be split into three parts
	\begin{enumerate}[(i)]
		\item State your inductive hypothesis $P(n).$

		\item State and justify the correctness of the base case.

		\item Prove the correctness of the induction step. \\
			This must explicitly state exactly in which step(s) the proof is using the
			inductive hypothesis.\\
			Also, every line of your proof should follow easily from the previous line.
			It's your job to provide justification of each step from the previous one,
			not the reader's.
	\end{enumerate}

	\subsection*{\textbf{Solution:}}
	\begin{enumerate}[(i)]
		\item Assume that the statement is true for some arbitrary natural $k$,
			meaning $\sum_{i=0}^{k}a_{i}=\frac{a_{0}-a_{k+1}}{1-r}$, where $a_{i}=a_{0}
			r^{i}$ and $r\neq1$.

		\item For the base case, let $n=0$. The left side is $\sum_{i=0}^{0}a_{i}=a_{0}$,
			and the right side is
			$\frac{a_{0}-a_{1}}{1-r}=\frac{a_{0}-a_{0}r}{1-r}=\frac{a_{0}(1-r)}{1-r}=a_{0}$.
			Since both sides are equal, the statement holds for $n=0$.

		\item For $k+1$:
			\[
				\begin{aligned}
					\sum_{i=0}^{k+1}a_{i} & =\sum_{i=0}^{k}a_{i}+a_{k+1}              & \text{(expand the summation)}             \\
					                      & =\frac{a_0-a_{k+1}}{1-r}+a_{k+1}          & \text{(inductive hypothesis)}             \\
					                      & =\frac{a_0-a_{k+1}+a_{k+1}(1-r)}{1-r}     & \text{(combine using common denominator)} \\
					                      & =\frac{a_0-a_{k+1}+a_{k+1}-a_{k+1}r}{1-r} & \text{(distribute)}                       \\
					                      & =\frac{a_0-a_{k+1}r}{1-r}                 & \text{(cancel like terms)}                \\
					                      & =\frac{a_0-a_{k+2}}{1-r}.                 & \text{(since }a_{k+2}=a_{k+1}r\text{)}
				\end{aligned}
			\]
	\end{enumerate}

	\newpage

	\section*{\textbf{EC: P3.4.6} [10 pts]}
	A \textbf{Fermat number} is a natural of the form $F_{i} = 2^{2^i}+ 1$, where
	$i$ is any natural. In 1730 Goldbach used Fermat numbers to give an alternate proof
	that there are infinitely many primes.
	\begin{enumerate}[(a)]
		\item List the Fermat numbers $F_{0}, F_{1}, F_{2}, F_{3}$, and $F_{4}$.

		\item Prove that for any $n$, the product
			$F_{0} \cdot F_{1} \cdot \ldots \cdot F_{n}$ is equal to $F_{n+1}- 2$.

			{\em Hint. we recommend proving this by induction. If you do prove it via induction, don't forget to split it into three parts, as in problems 7 and 8.}

		\item Argue that no two different Fermat numbers can share a prime factor.
			Since there are infinitely many Fermat numbers, there must thus be infinitely
			many primes.
	\end{enumerate}

	\subsection*{\textbf{Solution:}}
	\begin{enumerate}[(a)]
		\item The The first five Fermat numbers are:
			\[
				\begin{aligned}
					F_{0} & = 2^{2^0}+1=3      \\
					F_{1} & = 2^{2^1}+1=5      \\
					F_{2} & = 2^{2^2}+1=17     \\
					F_{3} & = 2^{2^3}+1=257    \\
					F_{4} & = 2^{2^4}+1=65537.
				\end{aligned}
			\]

		\item

		\item
	\end{enumerate}
\end{document}